\section{Conflict-Free Replicated Data Types}
\label{sec:conflict-free-replicated-data-types}

Alice's laptop shows \enquote{Review} and Bob's phone shows \enquote{Ready} for the same task.
Sending each device the other's title does not tell either device which one to display.
If each simply adopts the incoming title, Alice ends up with \enquote{Ready} and Bob with \enquote{Review}: the titles have traded places, but the replicas still disagree.
The two edits are \emph{concurrent} because each device saved its change before learning of the other's.

\emph{Conflict-free replicated data types} (CRDTs) allow replicas to accept updates independently.
Each type defines rules for applying local edits and incorporating remote changes, so replicas that have received the same updates agree on the result.
Given eventual delivery and termination, this provides the strong eventual consistency described in Section~\ref{sec:consistency-models}.
A server can carry the changes without deciding which edit to accept.
\parencite[1--2]{crdt-overview}

\subsection{How replicas synchronize}

In \emph{state-based replication}, an edit changes the state of one replica.
When that replica synchronizes, it sends a copy of its current state to a peer, which merges the received state with its own.
A single state can therefore convey the effects of several local edits.
\parencite[Section~2.1]{conflict-free-replicated-data-types}

In \emph{operation-based replication}, the replica where an edit originates prepares a message describing the update.
Each replica, including the originating replica, applies the update to its own state.
In the model of \textcite[Section~2.4]{conflict-free-replicated-data-types}, delivery is reliable and respects causal order: an update is delivered before any update that depends on it.
Concurrent updates may arrive in different orders, so applying them in either order must yield equivalent states.

Figure~\ref{fig:replication-models} compares the two approaches for a set that supports additions only.
Both replicas start empty, and replica A adds \(x\) and then \(y\).
On the left, B receives the set \(\{x,y\}\) and merges it by taking the union; on the right, B applies the two additions.

\begin{illustration}[htbp]
	\centering
	\caption{The same two additions under state-based and operation-based replication}
	\label{fig:replication-models}
	\begin{tikzpicture}[thesis]
	\foreach \offset/\heading in {0/State-based,7.2/Operation-based}{
		\begin{scope}[xshift=\offset cm]
			\node[thesis/heading] at (2.7,1) {\heading};
			\foreach \level/\replica in {0/A,-1.6/B}{
				\node[anchor=east] at (0,\level) {\replica};
				\draw[thesis/lifeline,->] (0,\level) -- (5.8,\level);
			}
			\foreach \pos/\itemname in {0.8/x,3/y}{
				\fill (\pos,0) circle (1.7pt);
				\node[above] at (\pos,0) {Add \(\itemname\)};
			}
		\end{scope}
	}
	\draw[thesis/divider] (6.4,1.15) -- (6.4,-2.25);
	\draw[thesis/network,->] (3.3,0) -- node[right] {State \(\{x,y\}\)} (4.7,-1.6);
	\node[below] at (4.7,-1.6) {Merge by union};
	\begin{scope}[xshift=7.2cm]
		\foreach \start/\endpos/\itemname in {0.8/2.2/x,3/4.4/y}{
			\draw[thesis/network,->] (\start,0) -- node[right] {Add \(\itemname\)} (\endpos,-1.6);
			\node[below] at (\endpos,-1.6) {Apply};
		}
	\end{scope}
\end{tikzpicture}

	\figuresource{Author's illustration adapted from \textcite[Section~2, Figures~1--2]{conflict-free-replicated-data-types}.}
\end{illustration}
\FloatBarrier

\newpage
\subsection{Registers}

A register stores a value that can be read or replaced.
A \emph{last-writer-wins register} (LWW register) also stores the timestamp of the write that assigned the value.
In the state-based design, merging keeps the value--timestamp pair with the greater timestamp, while reading returns only the value.
Timestamps must be unique and ordered consistently with causality: a write made after observing another must rank above it.
Taking the maximum gives the same result in any merge order, and receiving a state twice has no further effect.
\parencite[Section~3.2.1]{comprehensive-study-crdts}

A \emph{hybrid logical clock} (HLC) combines a time component derived from wall clocks with a logical counter.
Each replica incorporates the timestamps it receives into its clock.
If its wall clock lags behind an observed write, the logical counter lets the replica give its next write a greater timestamp without waiting for physical time to catch up.
\parencite[Section~3.3]{logical-physical-clocks}
To order writes from different replicas, the register compares the time component first, then the counter, and finally the unique replica identifier if both HLC components match.

\subsection{Sets}

A replicated set needs a membership rule when one replica adds an element while another removes it.
An \emph{observed-remove set} gives each addition a unique identifier.
A removal cancels only those additions of the element that the removing replica has observed.
The element remains present if at least one addition survives, so an addition concurrent with a removal wins.
A remove-wins set makes the opposite choice for concurrent operations on the same element.
Both rules permit replicas to agree; they give different answers to the membership query.
\parencites[Section~3.3.5]{comprehensive-study-crdts}[3--5]{crdt-overview}

\subsection{Sequences}

A replicated sequence must preserve an order as well as record which elements are present.
An index alone cannot identify an element across replicas because an insertion or deletion can shift its position.
The \emph{Replicated Growable Array} (RGA) gives each inserted element a stable identifier and names its predecessor by identifier.
Concurrent insertions after the same predecessor both survive, ordered by descending insertion timestamp.
Deletion hides an element from reads but retains its identifier as a \emph{tombstone}, so a concurrent insertion can still locate its predecessor.
\parencite[Section~3.5.1]{comprehensive-study-crdts}

In Figure~\ref{fig:rga-tombstone}, \(C\) was inserted after \(B\).
Deleting \(B\) changes the visible sequence without erasing the anchor for \(C\).
\begin{illustration}[htbp]
	\centering
	\caption{An RGA retains a deleted element as an insertion anchor}
	\label{fig:rga-tombstone}
	\begin{tikzpicture}[thesis, element/.style={thesis/box, minimum width=2.2cm, minimum height=0.8cm}]
	\node[anchor=east] at (-1.4,0) {Stored elements};
	\node[element] (a) at (0,0) {\(A\)};
	\node[element,dashed] (b) at (3.2,0) {\(B\) (deleted)};
	\node[element] (c) at (6.4,0) {\(C\)};
	\draw[thesis/local,->] (a) -- (b);
	\draw[thesis/local,->] (b) -- (c);
	\node[anchor=east] at (-1.4,-1) {Query result};
	\node at (0,-1) {\([A,C]\)};
\end{tikzpicture}

	\figuresource{Author's illustration based on \textcite[Section~3.5.1, Specification~19]{comprehensive-study-crdts}.}
\end{illustration}
\FloatBarrier

\newpage
\subsection{Resolving concurrent title edits}

In the opening example, suppose Alice's laptop saves \enquote{Review} with timestamp \((100,0,A)\) and Bob's phone saves \enquote{Ready} with \((101,0,B)\), as shown in Figure~\ref{fig:lww-title-merge}.

\begin{illustration}[htbp]
	\centering
	\caption{Concurrent title updates in a state-based LWW register}
	\label{fig:lww-title-merge}
	\begin{tikzpicture}[thesis, state/.style={thesis/box, minimum width=2.8cm, minimum height=0.7cm}]
	\node[thesis/heading] at (0,0.7) {Local update \(u\)};
	\node[thesis/heading] at (4.2,0.7) {Exchange states};
	\node[thesis/heading] at (8.4,0.7) {After merge \(m\)};
	\node[anchor=east] at (-1.65,0) {Alice};
	\node[anchor=east] at (-1.65,-1.5) {Bob};
	\node[state] (a) at (0,0) {Review, \((100,0,A)\)};
	\node[state] (b) at (0,-1.5) {Ready, \((101,0,B)\)};
	\node[state] (aa) at (8.4,0) {Ready, \((101,0,B)\)};
	\node[state] (bb) at (8.4,-1.5) {Ready, \((101,0,B)\)};
	\draw[thesis/local,->] (a.east) -- (aa.west);
	\draw[thesis/local,->] (b.east) -- (bb.west);
	\draw[thesis/network,->] (a.east) -- (bb.west);
	\draw[thesis/network,->] (b.east) -- (aa.west);
\end{tikzpicture}

	\figuresource{Author's illustration adapted from \textcite[Figure~7]{comprehensive-study-crdts}.}
\end{illustration}
\FloatBarrier

Bob's write has the greater timestamp, so both replicas display \enquote{Ready}.
This does not establish which user typed later in real time.
The register discards Alice's concurrent rename.
